\documentclass[11pt]{article}

\usepackage[spanish]{babel}
\usepackage{cv}

\addbibresource{cv.bib}

\author{Miguel Masó, PhD}

\begin{document}
\maketitle
\begin{tabularx}{\linewidth}{lCr}
\email{miguel.maso@upc.edu} &
+34 717 712 217 &
Ingeniero de Caminos, Col. Nº 36285 \\[-.2em]
Ciudad: Barcelona (08006) &
Fecha de nacim.: 9/4/1989 &
\href{https://www.linkedin.com/in/miguel-maso/}{\faLinkedin \ linkedin.com/in/miguel-maso} \\[-.2em]
\end{tabularx}


\section{Perfil profesional}
Ingeniero de Caminos y Doctor en Mecánica Computacional con un perfil híbrido que combina \textbf{simulación física de alto nivel, desarrollo de software y automatización BIM}. Domino desde el código base (C++, C\# / .NET, Python) hasta el despliegue de herramientas en producción: en SOCOTEC y como autónomo llevé soluciones propias desde la teoría numérica hasta su uso operativo por equipos de ingeniería. Capacidades que son los componentes técnicos sobre los que se asientan la digitalización de infraestructuras y los \textbf{gemelos digitales}.

Como Responsable de Innovación y Desarrollo en SOCOTEC lideré la automatización de procesos de diseño y la interoperabilidad directa entre modelos de cálculo y BIM, eliminando el trabajo manual y liberando horas de ingeniería. Mi base doctoral en métodos numéricos (FEM, CFD, interacción fluido-estructura y aguas someras) aporta una profundidad en simulación física directamente transferible a la modelización hidráulica y a la resiliencia de infraestructuras.


\section{Competencias clave}
\begin{description}
    \item[Digitalización de infraestructuras]
        Interoperabilidad y automatización BIM (Revit API, Robot Structural Analysis API) e integración directa cálculo--BIM, desplegada en herramientas de producción.
        Modelización física predictiva (FEM, CFD, aguas someras) como base de simulación.
        Análisis y explotación de datos de ingeniería.
    \item[Algoritmia, optimización y datos]
        Optimización mediante algoritmos genéticos y metaheurísticos.
        Calibración y ajuste de modelos constitutivos; estadística avanzada y análisis de componentes principales (PCA).
        Modelos de orden reducido (ROM), incluidos los basados en redes neuronales.
        Gestión masiva de datos para el rescate y aceleración de entregas de proyecto.
    \item[Desarrollo de software]
        C++ (solvers de alto rendimiento), C\# / .NET (Framework 4.8 y .NET 6+), Python y Julia.
        Control de versiones: Git, flujos CI/CD. Plataformas: Windows, Linux.
    \item[Idiomas] Español (nativo) - Catalán (nativo) - Inglés (C1) - Italiano (B1)
\end{description}


\section{Experiencia orientada a digitalización}

\subsection{Automatización e interoperabilidad cálculo--BIM --- SOCOTEC Spain (2022--2025)}
\parbox{\linewidth}{
    \emph{Ingeniero Especialista, Responsable de Innovación y Desarrollo} \newline
    \emph{Necesidad}: los procesos manuales entre modelos de cálculo y BIM generaban retrabajo constante, errores de transcripción y cuellos de botella en las entregas. \newline
    \emph{Solución}: lideré el desarrollo de soluciones propias (C\# / .NET / Python) para la interoperabilidad directa entre cálculo estructural y BIM (Revit/Robot), definiendo prioridades y coordinando al equipo técnico. \newline
    \emph{Impacto}: eliminación del retrabajo manual y errores de transcripción, liberando horas de ingeniería en cada proyecto; en situaciones críticas, la gestión masiva de datos permitió sacar adelante entregas comprometidas.}

\subsection{Aplicación de seguridad y automatización en izado --- Hormipresa (2024--2025)}
\parbox{\linewidth}{
    \emph{Consultor externo (autónomo)} \newline
    \emph{Necesidad}: garantizar la seguridad del izado de piezas prefabricadas. \newline
    \emph{Solución}: desarrollé, por cuenta propia, una aplicación para verificar la resistencia de los anclajes de izado, automatización del recuento de material y generación de planos. \newline
    \emph{Impacto}: refuerzo de la seguridad en obra, ahorro de horas a proyectistas al eliminar tareas repetitivas.}

\subsection{Modelización física para infraestructura crítica (base de simulación)}
\parbox{\linewidth}{
    Simulación física avanzada, base para la componente de modelización de sistemas digitales predictivos: aguas someras y flujos de superficie libre (fundamento físico y numérico de la modelización hidráulica y de inundaciones), CFD para cargas de viento (Torre Garellano, Arup) y simulaciones CFD-estructura termo-plástica frente a fuego (SOCOTEC Francia). Análisis dinámico y de robustez (EN 1991-1-7) en structura singular para Zaha Hadid (Mercury Tower, Malta) y SACYR (New Velindre Cancer Centre, Cardiff).}


\section{Experiencia profesional}
\begin{tabularx}{\linewidth}{lX}
    2025 - actual &
    \textbf{Centro Internacional de Métodos Numéricos en la Ingeniería (CIMNE)} \newline
    Investigador postdoctoral \\

    2022 - actual &
    \textbf{Universitat Politècnica de Catalunya (UPC)} - Escuela de Ingeniería de Caminos \newline
    Profesor asociado \\

    2022 - 2025 &
    \textbf{SOCOTEC Spain} - Edificación y ciudades \newline
    Ingeniero Especialista, Responsable de Innovación y Desarrollo \\

    2017 - 2022 &
    \textbf{Centro Internacional de Métodos Numéricos en la Ingeniería (CIMNE)} \newline
    Estudiante de doctorado (beca FPI) \\
\end{tabularx}


\section{Formación}
\begin{tabularx}{\linewidth}{lX}
    2017 - 2022 &
    \textbf{Doctorado en Ingeniería Civil (Mecánica Computacional)} \newline
    UPC, Barcelona, España \\

    2007 - 2013 &
    \textbf{Ingeniero de Caminos, Canales y Puertos} (equivalente a MSc) \newline
    UPC, Barcelona, España \\

    2013 - 2015 &
    \textbf{Grado en Filosofía y Teología} \newline
    Pontificia Università della Santa Croce, Roma, Italia \\
\end{tabularx}


\section{Producción científica y otras actividades}
\begin{itemize}
    \itemsep=-.3em
    \item \textbf{Publicaciones}: autor de artículos en revistas indexadas sobre modelización numérica (aguas someras, interacción fluido-estructura, elementos finitos). \href{https://orcid.org/0000-0003-0962-8550}{\faOrcid \ ORCID 0000-0003-0962-8550}.
    \item \textbf{Docencia universitaria} (2017 - actual): Profesor asociado en la UPC. Asignaturas: Dinámica de Estructuras (Máster) y Resistencia de materiales (Grado).
    \item \textbf{Estancia de investigación}: Investigador visitante en el CIMAT (México); cálculo en tiempo real; formulaciones estabilizadas para leyes de conservación hiperbólicas (2017).
    \item \textbf{Divulgación}: Ponente en la Asociación de Consultores de Estructuras (ACE) sobre sostenibilidad y descarbonización de estructuras (2024, 2025).
\end{itemize}


\end{document}
